\originalchapter{8.422公开作业五至十}
\originalsection{作业五: 散射与色散力}
\source{Assignment 5, 第1题.}
\begin{exercise}弹性散射$\ket{a;\mathbf k,\epsilon}\to\ket{a;\mathbf k',\epsilon'}$. (a) 分别用$-\mathbf d\cdot\mathbf E$和Coulomb规范$-q\mathbf p\cdot\mathbf A/m+q^2A^2/(2m)$算最低非零$T$矩阵, 比较低频Rayleigh与高频Thomson; (b) 证$\sigma_T=8\pi r_e^2/3$. 原题虽引用教材习题, 本身给了完整求解条件.\end{exercise}
\begin{solution}以电场模式归一化$\mathcal E_k=\sqrt{\hbar\omega/(2\epsilon_0V)}$, 长度表象的二阶矩阵为
\[T_{fi}=\mathcal E_{k'}\mathcal E_k\sum_n\left[\frac{(\epsilon'^*\cdot\mathbf d_{an})(\epsilon\cdot\mathbf d_{na})}{E_a+\hbar\omega-E_n}+\frac{(\epsilon\cdot\mathbf d_{an})(\epsilon'^*\cdot\mathbf d_{na})}{E_a-\hbar\omega-E_n}\right].\]
对各向同性态等于$-\mathcal E_{k'}\mathcal E_k\alpha(\omega)\epsilon'^*\cdot\epsilon$. 速度表象把两偶极因子换为$p$并乘$q^2/(m^2\omega^2)$, 还加$A^2$的一阶接触项. 用$p_{na}=\ii m\omega_{na}r_{na}$, 把$\omega_{na}^2/(\omega_{na}^2-\omega^2)=1+\omega^2/(\omega_{na}^2-\omega^2)$展开, 常数部分由TRK和规则恰与接触项结合, 剩余即长度结果. 这同时说明删$A^2$会破坏规范一致性.

低频$\alpha\simeq\alpha(0)$,$\sigma_R\propto\omega^4$. 高频$\alpha\simeq-e^2/(m\omega^2)$, 微分截面$r_e^2|\epsilon'^*\cdot\epsilon|^2$, 对两末偏振和方向积分$\int[1-|\epsilon\cdot\hat k'|^2]\dd\Omega=8\pi/3$, 即$\sigma_T=8\pi r_e^2/3$. 高频仍须低于相对论Compton区间, 否则此非相对论模型失效.
\end{solution}
\source{Assignment 5, 第2题.}
\begin{exercise}两基态原子非延迟偶极作用$V=[\mathbf d_a\cdot\mathbf d_b-3(\mathbf d_a\cdot\hat R)(\mathbf d_b\cdot\hat R)]/R^3$采用Gaussian单位. (a) 求首个非零能移; (b) 用方向振子强度$f_{ig}=2m\omega_{ig}|x_{ig}|^2/\hbar$表达, 消去交叉方向项; (c) 两原子各以一条$f\simeq1$主跃迁近似, 用静态极化率表达$C_6$.\end{exercise}
\begin{solution}(a) 一阶为零, 二阶$-C_6/R^6$,$C_6=\sum_{ij}|\langle ij|\mathbf d_a\cdot\mathbf d_b-3d_{az}d_{bz}|gg\rangle|^2/[\hbar(\omega_i+\omega_j)]$. (b) 各向同性基态按整套简并激发态求和, 不同Cartesian方向交叉项为零, 张量$\operatorname{diag}(1,1,-2)$平方迹为6. 于是$C_6=6\sum_{ij}e^4|x_i|^2|x_j|^2/[\hbar(\omega_i+\omega_j)]=3\hbar e^4\sum_{ij}f_if_j/[2m^2\omega_i\omega_j(\omega_i+\omega_j)]$. (c) $\alpha_a=2e^2|x_a|^2/(\hbar\omega_a)$, 得$C_6=3\hbar\alpha_a\alpha_b\omega_a\omega_b/[2(\omega_a+\omega_b)]$. SI形式再除$(4\pi\epsilon_0)^2$; 两相同原子为$3\hbar\omega_a\alpha_a^2/4$.
\end{solution}

\originalsection{作业六: 共振交换与相位噪声}
\source{Assignment 6, 第1题.}
\begin{exercise}两二级原子一激发一基态. (a) 跃迁能近乎相同但相差$\hbar\delta$, 大$R$两态如何移能, 何时微扰失效? (b) 完全同频的势曲线; (c) 同频时$R^{-3}$系数与自发率关系.\end{exercise}
\begin{solution}在$eg,ge$基中去掉平均能量,$H=\bigl(\begin{smallmatrix}\hbar\delta/2&C_3/R^3\\C_3^*/R^3&-\hbar\delta/2\end{smallmatrix}\bigr)$. 根$\pm\sqrt{(\hbar\delta/2)^2+|C_3|^2/R^6}$. (a) 弱交换的移能为$\pm|C_3|^2/(\hbar|\delta|R^6)$, $R\lesssim(|C_3|/\hbar|\delta|)^{1/3}$时简并处理不可省. (b) 完全共振分支$\pm|C_3|/R^3$, 态为对称反对称组合. (c) 对固定偶极方向$C_3=\mathcal Gd^2/(4\pi\epsilon_0)$,$\mathcal G=1-3\cos^2\vartheta$, 用式\eqref{eq:spont}得$C_3=3\mathcal G\hbar\Gamma/(4k^3)$. 角度与简并支不同会改变系数, 不能给无方向的单一常数.
\end{solution}
\source{Assignment 6, 第2题.}
\begin{exercise}把电子模型化为两块面积$a^2$、间隔$a$的平板, 各带同号电荷$e/2$. 平衡Casimir吸引与静电斥力, 求Gaussian单位$\alpha=e^2/(\hbar c)$.\end{exercise}
\begin{solution}表面密度$\sigma=e/(2a^2)$, 每板受另一板场$2\pi\sigma$, 斥力$\pi e^2/(2a^2)$. 理想无限板Casimir吸力为$\pi^2\hbar c/(240a^2)$. 平衡给$\alpha=\pi/120\simeq0.0262$, 约真实值的3.6倍. 面积边长与间距相同不满足无限大板条件, 模型也没有电子内部结构的实验支持; 这是题给玩具模型的结果, 不可称为推导出基本常数.
\end{solution}
\source{Assignment 6, 第3题.}
\begin{exercise}任意原子密度矩阵. (a) $R_z(\theta)$随机相位均值0方差$2\lambda t$, 求条件与平均态; (b) 环境两态初始0,$H_{SE}=\ket e\bra e\otimes(\gamma\ket0\bra1+\gamma^*\ket1\bra0)$, 求$a\ket g+b\ket e$的短时联合演化及原子态; (c) 以概率$(1-\e^{-\lambda t})/2$做$Z$翻转, 求平均态; 比较三模型是否完全等价.\end{exercise}
\begin{solution}(a) 对角元不变, 非对角乘$\e^{\pm\ii\theta}$, 平均乘$\e^{-\lambda t}$. (b) 令$\gamma$以角频率计, 联合态$a\ket{g,0}+b\ket e[\cos(|\gamma|t)\ket0-\ii\gamma^*\sin(|\gamma|t)\ket1/|\gamma|]$. 部分迹使非对角乘$\cos(|\gamma|t)=1-|\gamma|^2t^2/2+\cdots$, 布居不变. (c) $(1-p)\rho+pZ\rho Z$使非对角乘$1-2p=\e^{-\lambda t}$. (a,c)精确同信道; 固定单环境模式(b)只有短时平方退相干并会复现, 不等于全时段指数. 要让(b)成为相同Markov过程, 必须每步重置环境并取$|\gamma|^2\dd t/2=\lambda$的碰撞极限. 原件“完全等价”缺这个额外条件.
\end{solution}

\originalsection{作业七: Bloch、真空模与缀饰态}
\source{Assignment 7, 第1题.}
\begin{exercise}初态$\rho_{ee}=\rho_{ge}=0$, 弱激发$|\Omega|\ll\Gamma$. (a) 推导$\rho_{ee}=\frac{|\Omega|^2/4}{\delta^2+\Gamma^2/4}[1+\e^{-\Gamma t}-2\cos(\delta t)\e^{-\Gamma t/2}]$, 求$\Gamma\to0$极限; (b) 原件声称只要$|\Omega|t\ll1$便有$\rho_{ee}=|\Omega|^2t^2/4$, 独立于失谐与衰减, 检查所需条件并解释.\end{exercise}
\begin{solution}(a) 最低阶激发幅度$c_e=-\ii\Omega[1-\e^{-(\Gamma/2-\ii\delta)t}]/[2(\Gamma/2-\ii\delta)]$, 取模平方得式. 零线宽给$|\Omega|^2\sin^2(\delta t/2)/\delta^2$, 共振极限为平方律. (b) 短时展开还要求$|\delta|t\ll1,\Gamma t\ll1$; 仅驱动弱不够, 远失谐时即使$\Omega t$小仍可有振荡. 满足真正短时条件时能量尚未分辨, 同一初始耦合幅度决定平方律.
\end{solution}
\source{Assignment 7, 第2题.}
\begin{exercise}单模Jaynes--Cummings的零、一激发基$\ket{0g},\ket{0e},\ket{1g}$, 后两块$H/\hbar=(1/2)\bigl(\begin{smallmatrix}-\delta&\Omega_1\\\Omega_1&\delta\end{smallmatrix}\bigr)$, 零激发能$-\hbar(\delta/2+\omega_0)$. (a) 求完整$U$; (b) 初态$0e$, 测到0或1光子的原子态; (c) 初腔真空、任意原子混态的约化演化; (d) 每次短相互作用后测量并重置腔, 求重复信道固定点; (e) 把$e,g$画为Bloch南北极, 画$\ket\pm=(e\pm g)/\sqrt2$与$e$轨迹; (f) 为什么短时结果不依失谐, 如何与远失谐抑制发射协调?\end{exercise}
\begin{solution}记$Q=\sqrt{\delta^2+\Omega_1^2}$,$C=\cos(Qt/2),S=\sin(Qt/2)$. (a) 零激发块相位$z=\e^{\ii(\delta/2+\omega_0)t}$, 其余块为$\bigl(\begin{smallmatrix}C+\ii\delta S/Q&-\ii\Omega_1S/Q\\-\ii\Omega_1S/Q&C-\ii\delta S/Q\end{smallmatrix}\bigr)$. (b) 无光子时原子$e$, 一光子时$g$, 概率$|q|^2$和$1-|q|^2$,$q=C+\ii\delta S/Q$. (c) 在$g,e$基,$\rho_{ee}'=|q|^2\rho_{ee}$,$\rho_{gg}'=\rho_{gg}+(1-|q|^2)\rho_{ee}$,$\rho_{ge}'=zq^*\rho_{ge}$. (d) 单次损失概率$p=\Omega_1^2\Delta t^2/4+O(\Delta t^4)$, 重复$n$次布居乘$(1-p)^n$, 相干乘$(\sqrt{1-p})^n$及自由相位; $p>0$固定点为$g$. 固定耦合下$\Delta t\to0$的每秒率$p/\Delta t\to0$, 是Zeno极限; 要保持非零连续率需缩放耦合或真实连续模浴, 不能称固定单腔极限自动给有限自发率.

(e) 旋转系去相位,$x,y\mapsto\sqrt{1-p}(x,y)$,$z\mapsto(1-p)z+p$, 这里$g$北极. 三轨迹分别$(\pm\e^{-\gamma t/2},0,1-\e^{-\gamma t})$和$(0,0,1-2\e^{-\gamma t})$. 前两条在$x,z$截面满足$z=1-x^2$, 后一条沿轴, 可直接据此画而无假定三维透视. (f) 极短$t\ll1/|\delta|$时不能分辨失谐能量; 远失谐抑制是$t$能分辨$\delta$之后的结果. 把短时极限与大失谐极限交换会产生错误结论.
\end{solution}
\source{Assignment 7, 第3题.}
\begin{exercise}$H/\hbar=\omega_0Z/2+\Omega_1[X\cos\omega_Lt+Y\sin\omega_Lt]/2$. (a) 取$\psi=a\e^{\ii\omega_1t}g+b\e^{\ii\omega_2t}e$, 选频率消时间; (b,c) 写失谐Hamiltonian, 用$\sin2\theta=\Omega_1/Q$表达; (d) 求根和态; (e) 回到原表象.\end{exercise}
\begin{solution}选$\omega_1=\omega_L/2,\omega_2=-\omega_L/2$, $\delta_L=\omega_L-\omega_0$. 在$e,g$基$H'/\hbar=(-\delta_LZ+\Omega_1X)/2$, $Q=\sqrt{\delta_L^2+\Omega_1^2}$. 取$\cos2\theta=-\delta_L/Q$, 正支态$\cos\theta e+\sin\theta g$, 负支态$-\sin\theta e+\cos\theta g$, 能量$\pm\hbar Q/2$. 各支乘$\e^{\mp\ii Qt/2}$, 再分别给$e,g$乘$\e^{-\ii\omega_Lt/2},\e^{\ii\omega_Lt/2}$, 任意初态按这两支展开便得完整解.
\end{solution}

\originalsection{作业八: 光力与腔损耗}
\source{Assignment 8, 第1题.}
\begin{exercise}谐振束缚电子满足$m\ddot r+\gamma\dot r+m\omega_0^2r=-eE_0(x_0)\e^{\ii(\theta(x_0)+\omega t)}$, 核固定$x_0$. 偶极$p=\hat p[u\cos\Theta-v\sin\Theta]$. (a) 时均$(p\cdot\hat E)\nabla E$; (b) 用瞬时偶极势能重算并解释差别; (c) 解$p$并近共振写力; (d) 比较弱场Bloch光力, 用振子强度解释.\end{exercise}
\begin{solution}(a) 相同频率平均$\cos^2=\sin^2=1/2$, 交叉为0, 得$\overline F=(\hat p\cdot\hat E)[u\nabla E_0+vE_0\nabla\theta]/2$. (b) 对已给偶极在固定$p$下对$-pE$取梯度与(a)相同. 若错误地先代$p=\alpha E$再微分, 保守项会多一倍; 诱导偶极势应为$-\alpha E^2/2$, 因建立极化也需能量. 吸收造成的相位滞后项不能由静态保守势表达. (c) 复极化率$\alpha=e^2/[m(\omega_0^2-\omega^2+\ii\Gamma\omega)]$, $\Gamma=\gamma/m$. $u=\Re\alpha E_0,v=\Im\alpha E_0$, 近共振$\alpha=-e^2(2\delta+\ii\Gamma)/[m\omega_0(4\delta^2+\Gamma^2)]$. 因而
\[F=-\frac{e^2}{2m\omega_0}\frac{\delta\nabla E_0^2+\Gamma E_0^2\nabla\theta}{4\delta^2+\Gamma^2}.\]
(d) $\Omega_R^2=e^2|z_{fi}|^2E_0^2/\hbar^2$, $f=2m\omega_0|z_{fi}|^2/\hbar$, 故$\hbar\Omega_R^2=f e^2E_0^2/(2m\omega_0)$, 逐项给$F_Q=fF_C$. 偶极的实际强度由量子矩阵元取代经典单位振子.
\end{solution}
\source{Assignment 8, 第2题.}
\begin{exercise}腔$H_0=\hbar\omega a^\dagger a$, 短时跳跃概率$\Gamma\dd t\langle n\rangle$. 跳跃未归一态$\sqrt{\Gamma\dd t}a\ket\psi$, 无跳跃由$H_{
m eff}=H_0-\ii\hbar\Gamma n/2$演化. (a) 归一化两态与解释损失项; (b) 重组混合密度仅含原$\psi$; (c) 一阶展开推导Lindblad并找$C$.\end{exercise}
\begin{solution}$\ket{\psi_1}=a\ket\psi/\sqrt{\langle n\rangle}$,$\ket{\psi_0}=K_0\ket\psi/\sqrt{1-\Gamma\dd t\langle n\rangle}$, $K_0=I-\ii H_0\dd t/\hbar-\Gamma n\dd t/2$. $n$多者更易泄漏, 故无计数对高数态幅度压低更快. 两未归一态外积直接相加,$\rho'=K_0\rho K_0^\dagger+\Gamma\dd t\,a\rho a^\dagger$. 展开得$\dot\rho=-\ii[H_0,\rho]/\hbar+\Gamma[a\rho a^\dagger-\{n,\rho\}/2]$, $C=\sqrt\Gamma a$, $C^\dagger=\sqrt\Gamma a^\dagger$. 非Hermitian项只描述条件无跳跃, 完整未条件化信道仍保持迹.
\end{solution}

\originalsection{作业九: 减速与MOT密度}
\source{Assignment 9, 第1题.}
\begin{exercise}Na炉600 K, 自然角宽$\Gamma=2\pi\times10$ MHz, Zeeman斜率$2\pi\times1.4$ MHz/G,$\lambda=589$ nm. 原子沿$x$运动、光反向, 失谐$\delta$, 以$f a_{\max}$减速到零,$0<f<1$. (a) 最大减速; (b) $B(x)$与长度, 初速为热速率分布峰值; (c) 散射率$\Gamma_s$时, 发射沿$\pm x$、各向同性、偶极图案三种纵向动量扩散.\end{exercise}
\begin{solution}$a_{\max}=\hbar k\Gamma/(2m)\simeq9.3\times10^5$ m/s$^2$. 采用Maxwell速率峰$v_p=\sqrt{2k_BT/m}\simeq659$ m/s, $v(x)=\sqrt{v_p^2-2fa_{\max}x}$,$L=v_p^2/(2fa_{\max})\simeq0.234/f$ m. 共振$\delta+kv-g_BB=0$, 所以$B(x)=[\delta+kv(x)]/g_B$. 若采用出炉流量加权峰$\sqrt{3k_BT/m}$, 长度增为3/2, 原题未区分需说明. 定义$D$使$\dd\operatorname{Var}p_x/\dd t=2D$. 吸收事件计数Poisson贡献$(\hbar k)^2\Gamma_s$, 发射再贡献此量乘$\langle\cos^2\vartheta\rangle$. 三种$D$分别为$(\hbar k)^2\Gamma_s$, $2(\hbar k)^2\Gamma_s/3$, 对垂直于$x$的线偶极图案为$7(\hbar k)^2\Gamma_s/10$; 若偶极沿$x$, 为$3(\hbar k)^2\Gamma_s/5$. 题未定偶极方向, 不存在唯一第三值.
\end{solution}
\source{Assignment 9, 第2题.}
\begin{exercise}同样Na初速$v_p$, 反向激光对静止原子共振,$I=5I_s$. 求减到零的时间与路程.\end{exercise}
\begin{solution}$-\dot v=a_{\max}s/[1+s+4k^2v^2/\Gamma^2]$, $s=5$. 分别积分$\dd t$及$v\dd t$, 得
\[t=\frac{(1+s)v_p+4k^2v_p^3/(3\Gamma^2)}{a_{\max}s},\qquad L=\frac{(1+s)v_p^2/2+k^2v_p^4/\Gamma^2}{a_{\max}s}.\]
代题值约$t=2.4$ s,$L=1.18\times10^3$ m, 显著长于Zeeman共振减速. 速度零是此积分的终止点; 同一激光若继续开, 会把原子推向反向而非自动停住.
\end{solution}
\source{Assignment 9, 第3题.}
\begin{exercise}三维MOT六束各强度$I$, 原子吸收泵浦截面$\sigma_L$, 吸收荧光截面$\sigma_R$, 密度$n$. (a) 从两原子间反平方斥力证$\nabla\cdot F_R=6\sigma_L\sigma_RIn/c$; (b) 解释衰减吸引项$\nabla\cdot F_A=-6\sigma_L^2In/c$; (c) 加$F_T=-\kappa r$求密度上限; (d) 红失谐强驱动时比较$\sigma_R,\sigma_L$; (e) 仅比例$f$原子亮态时各力如何缩放和$n_{\max}$变化.\end{exercise}
\begin{solution}(a) 一个原子散射六束功率$6\sigma_LI$, 距$r$处荧光动量流给$F=6\sigma_L\sigma_RI/(4\pi cr^2)$. Gauss通量每源为$6\sigma_L\sigma_RI/c$, 乘密度即得式. (b) 远侧光经原子云衰减, 两束不平衡向内, 吸收和遮挡均用$\sigma_L$, 符号相反. (c) $-3\kappa+6I\sigma_L(\sigma_R-\sigma_L)n/c=0$, 所以$n_{\max}=\kappa c/[2I\sigma_L(\sigma_R-\sigma_L)]$, 要$\sigma_R>\sigma_L$. 若两者相等或更小, 本机制不产生有限上限. (d) 红失谐强光饱和使泵浦吸收变弱, 荧光谱有靠近原子共振的成分, 往往$\sigma_R>\sigma_L$, 具体值须光谱加权而非普遍恒等. (e) 若亮态比例均匀并快混合, 散射源和受力接收者均乘$f$, 两密度力平均乘$f^2$, 阱力乘$f$, 所以$n_{\max}$约增$1/f$. 暗核实际空间分布不均, 简单比例只作模型说明.
\end{solution}

\originalsection{作业十: 缀饰态、边带与能量输运}
\source{Assignment 10, 第1题.}
\begin{exercise}沿作业七的经典缀饰态: (a) 弱耦合$\Omega_1\ll|\delta_L|$的AC移位和态; (b) 用缀饰偶极平均求动态极化率, 代振子强度与$e^2f/[m(\omega_0^2-\omega_L^2)]$比较, 讨论共振; (c) 比较静态极限与近似缺项; (d) 裸态初始展开缀饰态求原态概率; (e) 两支的受激力; (f) 经典时变与量子场定态两种图像的联系.\end{exercise}
\begin{solution}(a) 基态光移$\hbar\Omega_1^2/(4\delta_L)$, 激发态相反, 混合幅度约$\Omega_1/(2\delta_L)$. (b) 平均偶极与$E_0$同频, 旋波模型$\alpha_{\rm RWA}=d^2/[\hbar(\omega_0-\omega_L)]$; 全式$\alpha=2d^2\omega_0/[\hbar(\omega_0^2-\omega_L^2)]$, 使用方向$f=2m\omega_0d^2/(\hbar e^2)$即题给经典形式. 共振处无阻尼弱场表达发散, 必须计饱和和衰减. (c) $\omega_L\to0$全式为$2d^2/(\hbar\omega_0)$, RWA为一半, 缺反旋项. (d) 裸态以两缀饰态展开, 演化相位差$Qt$, 原态概率$1-\Omega_1^2\sin^2(Qt/2)/Q^2$. (e) 缀饰能$E_\pm=\pm\hbar Q/2$, 力$F_\pm=\mp\hbar\Omega_1\nabla\Omega_1/(2Q)$, 若失谐也有空间变化需加其梯度. (f) 量子场的$\ket{e,n},\ket{g,n+1}$为时间独立联合态; 大相干场窄相对光子数涨落时$g\sqrt{n+1}$趋经典耦合, 旋转系除去载频, 两描述一致. 有限光子数时场原子纠缠与collapse/revival不能由固定经典场完全表达.
\end{solution}
\source{Assignment 10, 第2题.}
\begin{exercise}离子运动角频$\nu$, $\eta=k\sqrt{\hbar/(2m\nu)}\ll1$, 激光$H_I=\hbar\Omega(\sigma_++\sigma_-)\cos(kx-\omega t)$. (a) 展开并求红蓝边带Rabi频率; (b) 低饱和$\Omega\ll\Gamma$写$\dot p_n=nA_+p_{n-1}+(n+1)A_-p_{n+1}-[(n+1)A_++nA_-]p_n$, 包括吸收边带与自发反冲两路径, 求$A_\pm$; (c) 解平均声子数和分辨边带红失谐$\delta=-\nu$的温度与时间; (d) $^{88}$Sr$^+$674 nm,$\nu=2\pi\times10$ MHz,$\Gamma=2\pi\times0.4$ Hz,$\Omega=2\pi\times0.04$ Hz, 再改$\Gamma=2\pi\times100$ kHz,$\Omega=2\pi\times10$ kHz, 比较冷却与每秒数千声子的电场加热.\end{exercise}
\begin{solution}(a) 展开$\e^{\ii kx}=1+\ii\eta(a+a^\dagger)$, 共振慢项给$\Omega_{n,n-1}=\eta\Omega\sqrt n$, $\Omega_{n,n+1}=\eta\Omega\sqrt{n+1}$. (b) 令$\zeta$为发射反冲在该轴的几何平均因子, $L(u)=[u^2+\Gamma^2/4]^{-1}$. 两路径给
\[A_\pm=\frac{\eta^2\Omega^2\Gamma}{4}\left[L(\delta\mp\nu)+\zeta L(\delta)\right].\]
$A_-$用$\delta+\nu$, $A_+$用$\delta-\nu$. (c) 求和$\sum n\dot p_n$得$\dot{\bar n}=-(A_--A_+)\bar n+A_+$, 所以$\bar n(t)=\bar n_\infty+(\bar n_0-\bar n_\infty)\e^{-(A_--A_+)t}$, $\bar n_\infty=A_+/(A_--A_+)$. 红边带分辨极限$\bar n_\infty\simeq(\Gamma/4\nu)^2+\zeta(\Gamma/2\nu)^2$, $\tau_c\simeq\Gamma/(\eta^2\Omega^2)$. 用热分布等效温度$T=\hbar\nu/[k_B\ln(1+1/\bar n_\infty)]$, 不把低占据量子振子强行用$k_BT=\hbar\nu\bar n$.

(d) $\eta^2\simeq5.0\times10^{-4}$. 取题未指定的发射几何$\zeta=1/3$, 原窄线$\bar n\simeq2.3\times10^{-16}$, $T\simeq13\,\mu$K,$\tau_c\simeq8.0\times10^4$ s; 有加热时完全不实用. 展宽后$\bar n\simeq1.46\times10^{-5}$,$T\simeq43\,\mu$K,$\tau_c\simeq0.32$ s. 此参数的净冷却率仅约3.1 s$^{-1}$, 每秒数千声子加热仍会使$\bar n$很大; “线宽大于加热率”只是必要尺度比较, 还要提高实际冷却率或降低噪声. 原件未给发射几何, 数值温度依$\zeta$而变.
\end{solution}
\source{Assignment 10, 第3题.}
\begin{exercise}原子自库A经细管进入均匀远红失谐光区B, 光区长$d$、截面$A$、体积$V=Ad$. 原子密度在$\Delta t$内以$\dot n$线性增加. (a) 折射率$1+n\alpha/(2\epsilon_0)$, 求出口光频下降; (b) 积分光能减少; (c) 由Stark作用求原子能移并比较两能量, 解释动能来源. 无自发散射.\end{exercise}
\begin{solution}(a) 采用题的场相位$\omega t-kz$, 原子相移$-kdn\alpha/(2\epsilon_0)$, 因而$\Delta\omega=\omega d\alpha\dot n/(2\epsilon_0c)$. (b) 光子流$IA/(\hbar\omega)$, 总光能损失$\Delta E_{EM}=IA\Delta\omega\Delta t/\omega=IV\alpha\Delta n/(2\epsilon_0c)$. (c) 每原子Stark势$U=-\alpha I/(2\epsilon_0c)$, 共$V\Delta n$个新原子, 总势变化$-IV\alpha\Delta n/(2\epsilon_0c)$, 与光能损失数值相等、符号相反. 但此有效势是包含相互作用后的慢运动能, 不能再与裸场能按独立项重复相加. 从偶极模型看, 裸极化内能为$+\alpha E_0^2/4$、交互作用为$-\alpha E_0^2/2$, 净Stark能为$-\alpha E_0^2/4$; 进阱获得动能$+\alpha E_0^2/4$. 这些内能、动能与交互能共同变化, 连同光的红移满足总能量守恒. 未自发发射不表示场完全不输运能量; 受激过程与相移也携带能量交换.
\end{solution}
